% الجزء: sets-functions-relations
% الفصل: sets
% القسم: basics

\documentclass[../../../include/open-logic-section]{subfiles}

\begin{document}

\olfileid[ar]{sfr}{set}{bas}
\olsection{الامتدادية}

\emph{المجموعة} تجمّع من أشياء يُنظر إليه {بوصفه} شيئًا واحدًا.
وتُسمى الأشياء التي تتألف منها المجموعة \emph{عناصرَ} المجموعة أو
\emph{أعضاءَها}. فإذا كان $x$ هو !!a{element} من مجموعة~$A$، نكتب
$x \in A$؛ وإلا فنكتب $x \notin A$. وتُسمى المجموعة التي لا تحتوي على أي
!!{element}s المجموعة \emph{الخالية}، ويرمز
إليها بـ«$\emptyset$».

\begin{explain}
لا يهم كيف \emph{نعيّن} المجموعة، ولا كيف
\emph{نرتب} !!{element}sها، ولا حتى كم \emph{مرة} نعدّ
!!{element}sها. فالمهم وحده هو ماهية !!{element}sها.
ونصوغ ذلك في المبدأ الآتي.
\end{explain}

\begin{defn}[الامتدادية]
  إذا كانت $A$ و$B$ مجموعتين، فإن $A = B$ إذا وفقط إذا كان
  كل !!{element} من~$A$ أيضًا !!a{element} من~$B$، والعكس
  صحيح.
\end{defn}

تسوّغ الامتدادية بعض الرموز. وبوجه عام، عندما تكون لدينا بعض
الأشياء $a_{1}$، \dots، $a_{n}$، فإن $\{a_{1}, \dots, a_{n}\}$ هي
\emph{المجموعة الوحيدة} التي !!{element}sها هي $a_1, \ldots, a_n$. ونشدد
على وصفها بـ«\emph{الوحيدة}»، إذ تقضي الامتدادية بأنه لا يمكن
أن توجد إلا \emph{مجموعة واحدة} من هذا القبيل. بل إن الامتدادية تسوّغ
أيضًا ما يأتي:
  \[
    \{a, a, b\} = \{a, b\} = \{b,a\}.
  \]
وهكذا يتحقق المقصود: فعندما ننظر في المجموعات، لا نهتم
بترتيب !!{element}sها، ولا بعدد المرات التي
تُذكر فيها.

\begin{tagblock}{novice}
\begin{ex}
كلما كانت لدينا طائفة من الأشياء، أمكننا جمعها معًا في
مجموعة. فمجموعة إخوة ريتشارد وأخواته، مثلًا، مجموعة
تضم شخصًا واحدًا، ويمكن أن نكتبها $S=\{\textrm{روث}\}$.
ومجموعة الأعداد الصحيحة الموجبة الأصغر من $4$ هي $\{1, 2, 3\}$، لكنها
تُكتب أيضًا $\{3, 2, 1\}$، أو حتى $\{1, 2, 1, 2, 3\}$.
وهذه كلها المجموعة نفسها بموجب الامتدادية، لأن كل !!{element}
في $\{1, 2, 3\}$ هو أيضًا !!a{element} في $\{3, 2, 1\}$ (وفي $\{1,
2, 1, 2, 3\}$)، والعكس صحيح.
\end{ex}
\end{tagblock}

كثيرًا ما نعيّن مجموعة بواسطة خاصية تشترك فيها !!{element}sها.
وسنستعمل لذلك الترميز المختصر الآتي:
$\Setabs{x}{\phi(x)}$، حيث ترمز $\phi(x)$ إلى الخاصية
التي يجب أن يتصف بها~$x$ لكي يُعدّ من !!{element}s
المجموعة.

\begin{tagblock}{novice}
\begin{ex}
في مثالنا، كان يمكننا أيضًا تعيين $S$ كما يأتي:
\[
S = \Setabs{x}{x \text{ من إخوة ريتشارد وأخواته}}.
\]
\end{ex}
\end{tagblock}

\begin{tagblock}{math}
\begin{ex}
يُسمى العدد \emph{تامًا} إذا وفقط إذا كان مساويًا لمجموع
قواسمه الفعلية (أي الأعداد التي تقسمه من دون باق، ولا
تساوي العدد نفسه). فمثلًا، العدد $6$ تام لأن
قواسمه الفعلية هي $1$ و$2$ و~$3$، ولأن $6 = 1 + 2 + 3$. والواقع أن
$6$ هو العدد الصحيح الموجب الوحيد الأصغر من $10$ الذي يكون تامًا. ولذلك
يمكننا، باستخدام الامتدادية، القول:
\[
	\{6\} = \Setabs{x}{x\text{ تام و }0 \leq x \leq 10}
\]
ونقرأ الترميز في الطرف الأيمن هكذا: «مجموعة قيم $x$ التي يكون فيها $x$
تامًا و$0 \leq x \leq 10$». وتؤكد هذه المساواة
أننا، عندما ننظر في المجموعات، لا نهتم بكيفية تعيينها.
وعلى نحو أعم، تضمن الامتدادية أن هناك دائمًا
مجموعة واحدة فقط من قيم $x$ التي تحقق $\phi(x)$.
ومن ثم تسوّغ الامتدادية تسمية
$\Setabs{x}{\phi(x)}$ \emph{المجموعة الوحيدة} لقيم $x$ التي تحقق~$\phi(x)$.
\end{ex}
\end{tagblock}

تتيح لنا الامتدادية طريقة لإثبات تطابق المجموعات: فلإثبات
$A = B$، أثبت أنه كلما كان $x \in A$ كان أيضًا $x \in B$،
وكلما كان $y \in B$ كان أيضًا $y \in A$.

\begin{prob}
أثبت أنه توجد على الأكثر مجموعة خالية واحدة؛ أي أثبت أنه إذا كانت $A$ و$B$
مجموعتين بلا !!{element}s، فإن $A = B$.
\end{prob}

\end{document}
